\documentclass[10pt]{mypackage}

\usepackage[uprightgreeks,light]{kpfonts}
\renewcommand{\coloneq}{\coloneqq}
%\renewcommand{\mathbb}{\mathds}
\usepackage{homework}
%\usepackage{notes}

%\usepackage[ backend=bibtex, style = alphabetic, sorting=ynt ]{biblatex}
%\addbibresource{  }

\usepackage{parskip}

\fancyhf{}
\fancyhead[R]{Avinash Iyer}
\fancyhead[L]{Algebra II: Homework 9}
\fancyfoot[C]{\thepage}

\setcounter{secnumdepth}{0}

\begin{document}
\RaggedRight
\begin{problem}
  Let $\sigma$ be an action of a group $G$ on a finite set $X$, and let $\left( \C[X] \right)_{\sigma}$ be the corresponding permutation representation of $G$ over $\C$. The representation $\C[X]$ decomposes as a direct sum of two subrepresentations,
  \begin{align*}
    \left( \C[X] \right)_{\sigma} &= W_{\sigma} \oplus W_{\sigma}^{\perp},
  \end{align*}
  where
  \begin{align*}
    W_{\sigma} &= \set{\lambda \sum_{x\in X}x | \lambda\in \C}
  \end{align*}
  and $W_{\sigma}^{\perp}$ is the orthogonal complement with respect to the Hermitian inner product on $\C[X]$, where we naturally identify $\C[X]\cong \C^{\left\vert X \right\vert}$. That is,
  \begin{align*}
    W_{\sigma}^{\perp} &= \set{\sum_{x\in X}\lambda_x x | \sum_{x\in X}\lambda_x = 0}.
  \end{align*}
  Denote by $\sigma^2$ the diagonal action of $G$ on $X\times X$ given by $g\cdot \left( x,y \right) = \left( g\cdot x,g\cdot y \right)$. Observe that $\sigma^2$ stabilizes the diagonal subset $\Delta(X)\subseteq X\times X$ given by
  \begin{align*}
    \Delta(X) &= \set{\left( x,x \right) | x\in X}.
  \end{align*}
  Hence, the diagonal action also stabilizes $\left( X\times X \right)\setminus \Delta(X)$. Denote by $\sigma(2)$ the restriction of $\sigma^2$ to $\left( X\times X \right)\setminus \Delta(X)$.

  The \textit{rank} of $\sigma$ is defined to be the number of orbits of $\sigma(2)$, denoted $\operatorname{rk}(\sigma)$. We have that $\operatorname{rk}\left( \sigma \right) = 1$ if and only if $\sigma$ is $2$-transitive; that is, if $G$ acts transitively on ordered pairs of distinct elements of $X$.
  \begin{enumerate}[(a)]
    \item Let $\tau$ be the action of a group on a set $X$, and let $F$ be a field. Prove that $\Dim\left( \left( F[X] \right)_{\tau}^{G} \right)$ is equal to the number of orbits of $\tau$.
    \item Now, assume that $\sigma$ is the action of a finite group on a finite set $X$, and let $\chi_{\omega}$ be the character of $\left( \C[X] \right)_{\sigma}$. Prove that
      \begin{align*}
        \iprod{\chi_{\sigma}}{\chi_{\sigma}} &= \frac{1}{\left\vert G \right\vert}\sum_{g\in G} \chi_{\sigma^2}\left( g \right).
      \end{align*}
    \item Applying (a) to $\tau = \sigma^2$ and $F = \C$, along with a suitable result from Lecture 15, conclude that $ \iprod{\chi_{\sigma}}{\chi_{\sigma}} $ is the number of orbits of $\sigma^2$.
    \item Use (c) to prove that $W_{\sigma}^{\perp}$ is irreducible if and only if $\left\vert X \right\vert = 2$ or $\operatorname{rk}\left( \sigma \right) = 1$.
  \end{enumerate}
\end{problem}
\begin{solution}\hfill
  \begin{enumerate}[(a)]
    \item Let
      \begin{align*}
        u &= \sum_{x\in X} \lambda_x x
      \end{align*}
      be some element with $\lambda_x\in F$. Then, we have that $u\in \left( F[X] \right)_{\tau}^{G}$ if and only if, for all $g\in G$, we have
      \begin{align*}
        u &= \tau(g)(u)\\
          &= \tau(g)\left( \sum_{x\in X}\lambda_x x \right)\\
          &= \sum_{x\in X}\lambda_x \tau\left( g \right)(x)\\
          &= \sum_{x\in X} \lambda_{\tau\left( g^{-1} \right)(x)}x,
      \end{align*}
      whence $\lambda_{\tau\left( g^{-1} \right)(x)} = \lambda_x$ for all $g\in G$, In particular, this means that $\lambda_x = \lambda_{\tau\left( h \right)(x)}$ for all $h\in G$, or that the coefficients of $u$ are constant on orbits. Therefore, if we let $O_1,\dots,O_r$ be the orbits of $\tau\colon G\rightarrow \Sym(X)$, then 
      \begin{align*}
        u_i &= \sum_{x\in O_i} x
      \end{align*}
      forms a basis for $\left( F[X] \right)_{\tau}^{G}$, as $u\in \left( F[X] \right)_{\tau}^{G}$ if and only if its coefficients are constant on orbits, whence $\Dim\left( \left( F[X] \right)_{\tau}^{G} \right)$ is equal to the number of orbits.
    \item We express
      \begin{align*}
        \iprod{\chi_{\sigma}}{\chi_{\sigma}} &= \frac{1}{\left\vert G \right\vert}\sum_{g\in G} \chi_{\sigma}\left( g \right) \overline{\chi_{\sigma}\left(g\right)}.
      \end{align*}
      Since $\sigma$ is a permutation representation, we have that
      \begin{align*}
        \chi_{\sigma}\left( g \right) &= \left\vert \operatorname{fix}\left( g \right) \right\vert,
      \end{align*}
      where $\operatorname{fix}\left( g \right)$ denotes the fixed points of $g\in G$. Let $\set{x_1,\dots,x_r}$ be the fixed points of some element $g\in G$, meaning that $\chi_{\sigma}\left( g \right) = r$. Extending this action to $X\times X$, we have that $\left( x,y \right)\in \operatorname{fix}\left( g \right)$ if and only if
      \begin{align*}
        g\cdot \left( x,y \right) &= \left( g\cdot x,g\cdot y \right)\\
                                  &= \left( x,y \right),
      \end{align*}
      meaning that $x$ and $y$ are both in $\operatorname{fix}\left( g \right)$. In particular, this means that 
      \begin{align*}
        \operatorname{fix}_{\sigma^2}(g) &= \set{x_1,\dots,x_r}\times \set{x_1,\dots,x_r},
      \end{align*}
      whence there are $r^2$ such elements.
      \begin{align*}
        \chi_{\sigma^2}\left( g \right) &= \left\vert \operatorname{fix}_{\sigma^2}\left( g \right) \right\vert\\
                                        &= \left\vert \operatorname{fix}\left( g \right) \right\vert^2
      \end{align*}
      In particular, this gives that
      \begin{align*}
        \chi_{\sigma}\left( g \right) \overline{\chi_{\sigma}\left( g \right)} &= \left\vert \operatorname{fix}_{\sigma}\left( g \right) \right\vert^2\\
                                                                               &= \left\vert \operatorname{fix}_{\sigma^2}\left( g \right) \right\vert\\
                                                                               &= \chi_{\sigma^2}\left( g \right).
      \end{align*}
    \item We have
      \begin{align*}
        \iprod{\chi_{\sigma}}{\chi_{\sigma}} &= \frac{1}{\left\vert G \right\vert}\sum_{g\in G} \chi_{\sigma^2}\left( g \right)\\
                                             &= \dim\left( \left( F[X] \right)_{\sigma^2}^{G} \right)\\
                                             &= \left\vert \set{O_1,\dots,O_r} \right\vert.
      \end{align*}
    \item If $\left\vert X \right\vert = 2$, then $W_{\sigma}^{\perp}$ has dimension $1$, so it is irreducible by definition. Else, if $\operatorname{rk}\left( \sigma \right) = 1$, then $\sigma$ is $2$-transitive. The action of $\sigma\left( g \right)$ on $\C[X]$ can be expressed on the basis
      \begin{align*}
        \beta &= \set{x_1 + \cdots + x_n,x_1-x_2,x_1-x_3,\dots,x_1-x_n}
      \end{align*}
      via the block diagonal matrix
      \begin{align*}
        \left[ \sigma\left( g \right) \right]_{\beta} &= \begin{pmatrix}1 & 0 \\ 0 & B\end{pmatrix}
      \end{align*}
      for some $\left( n-1 \right)\times \left( n-1 \right)$ matrix $B$. In particular, this means 
      \begin{align*}
        \chi_{\sigma}(g) = 1 + \chi_{\sigma}\left( g \right)|_{W_{\sigma}^{\perp}}.
      \end{align*}
      Since $\sigma$ is $2$-transitive, it follows that
      \begin{align*}
        \iprod{\chi_{\sigma}}{\chi_{\sigma}} &= 2,
      \end{align*}
      meaning that the restriction to $W_{\sigma}^{\perp}$ has
      \begin{align*}
        \iprod{\chi_{\sigma}|_{W_{\sigma}^{\perp}}}{\chi_{\sigma}|_{W_{\sigma}^{\perp}}} &= 1,
      \end{align*}
      so the representation is irreducible.

      Now, suppose $W_{\sigma}^{\perp}$ is irreducible, so that
      \begin{align*}
        \iprod{\chi_{\sigma}|_{W_{\sigma}^{\perp}}}{\chi_{\sigma}|_{W_{\sigma}^{\perp}}} &= 1.
      \end{align*}
      If $\left\vert X \right\vert\neq 2$, then we may view the action of $\sigma$ on the space $\C\left[ X \right]$ with basis
      \begin{align*}
        \beta &= \set{x_1 + \cdots + x_n,x_1-x_2,\dots,x_1-x_n},
      \end{align*}
      to be in the block diagonal form
      \begin{align*}
        \left[ \sigma(g) \right]_{\beta} &= \begin{pmatrix}1 & 0 \\ 0 & B\end{pmatrix},
      \end{align*}
      with $B$ having dimension strictly greater than $1$. It follows then that
      \begin{align*}
        \iprod{\chi_{\sigma}}{\chi_{\sigma}} &= 2,
      \end{align*}
      so that $\operatorname{rk}\left( \sigma \right) = 1$.
  \end{enumerate}
\end{solution}
\begin{problem}[Problem 2]
  Let $p$ be prime, and let $G$ be the set of all functions from $\F_p$ to $\F_p$ that have the form $x\mapsto ax + b$ for some $a\in \F_p^{\times}$ and $b\in \F_p$.
  \begin{enumerate}[(a)]
    \item Prove that $G$ is a group with isomorphic to the group of matrices
      \begin{align*}
        H &= \set{ \begin{pmatrix}a & b \\ 0 & 1\end{pmatrix} | a\in \F_p^{\times},b\in \F_p }.
      \end{align*}
    \item The group $G$ has a natural action on $\F_p$ given by $g\cdot x = g(x)$ for all $g\in G$ and $x\in \F_p$. Prove that this action is $2$-transitive.
  \end{enumerate}
\end{problem}
\begin{solution}\hfill
  \begin{enumerate}[(a)]
    \item We observe that $G$ and $H$ have the same order of $p\left( p-1 \right)$, and that $G$ is a group since any non-constant linear map from $\F_p$ to $\F_p$ is a bijection. Therefore, if we show that there is an injective homomorphism from $G$ to $H$, then this homomorphism is automatically surjective.

      Consider the map from $G$ to $H$ given by 
      \begin{align*}
        ax + b &\mapsto \begin{pmatrix}a & b \\ 0 & 1\end{pmatrix}.
      \end{align*}
      If $a_1,a_2\in \F_{p}^{\times}$ and $b_1,b_2\in \F_p$, then we have
      \begin{align*}
        a_1\left( a_2x + b_2 \right) + b_1 &= a_1a_2 x + \left( a_1b_2 + b_1 \right).
      \end{align*}
      We have that the left hand side maps to
      \begin{align*}
        \begin{pmatrix}a_1 & b_1 \\ 0 & 1\end{pmatrix}\begin{pmatrix}a_2 & b_2 \\ 0 & 1\end{pmatrix} &= \begin{pmatrix}a_1a_2 & a_1b_2 + b_1\\0 & 1\end{pmatrix},
      \end{align*}
      which gives that this map is a homomorphism.

      To show that it is injective, observe that the identity on $H$ is given by
      \begin{align*}
        \begin{pmatrix}1 & 0 \\ 0 & 1\end{pmatrix},
      \end{align*}
      which is the image under $x\mapsto x$, which means the homomorphism $G\rightarrow H$ is injective.
    \item Given $\left( x_1,y_1 \right)$ and $\left( x_2,y_2 \right)$ with $x_i\neq y_i$ and $\left( x_1,y_1 \right)\neq \left( x_2,y_2 \right)$, we may find an affine transformation via the classic point-slope formula,
      \begin{align*}
        g(x) &= \left( y_2-y_1 \right)\left( x_2-x_1 \right)^{-1}\left( x-x_1 \right) + y_1.
      \end{align*}
      Since these are distinct points with $x_1\neq x_2$, it follows that the action is $2$-transitive.
  \end{enumerate}
\end{solution}
\begin{problem}[Problem 3]\hfill
  \begin{enumerate}[(a)]
    \item Let $K = \Q\left( \sqrt{2},\sqrt{3} \right)$. Prove that $\left[ K:\Q \right] = 4$.
    \item Let $L = \Q\left( \sqrt{2} + \sqrt{3} \right)$. Prove that $L = K$, and hence $\left[ L:\Q \right] = 4$ by (a).
  \end{enumerate}
\end{problem}
\begin{solution}\hfill
  \begin{enumerate}[(a)]
    \item We have the tower of field extensions $\Q\subseteq \Q\left( \sqrt{2} \right)\subseteq \Q\left( \sqrt{2},\sqrt{3} \right)$, and we observe that the minimal polynomial for $\sqrt{3}$ is still given by $x^2 - 3$ in $\Q\left( \sqrt{2} \right)$, whence
      \begin{align*}
        \left[ K:\Q \right] &= \left[ K:\Q\left( \sqrt{2} \right) \right]\left[ \Q\left( \sqrt{2} \right):\Q \right]\\
                            &=4.
      \end{align*}
    \item We immediately see that $L\subseteq K$, as $\sqrt{2} + \sqrt{3}\in \Q\left( \sqrt{2},\sqrt{3} \right)$. Working over $\C$, we observe that
      \begin{align*}
        \frac{1}{\sqrt{3} + \sqrt{2}} &= \sqrt{3} - \sqrt{2},
      \end{align*}
      meaning that this expression is contained in $L$, whence
      \begin{align*}
        \sqrt{2} &= \frac{1}{2}\left( \sqrt{2} + \sqrt{3} - \frac{1}{\sqrt{3} + \sqrt{2}} \right)\\
        \sqrt{3} &= \frac{1}{2} \left( \sqrt{2} + \sqrt{3} + \frac{1}{\sqrt{3} + \sqrt{2}} \right)
      \end{align*}
      are expressions wholly in $L$ for $\sqrt{2}$ and $\sqrt{3}$, meaning that $\sqrt{2}$ and $\sqrt{3}$ are contained in $L$, so $\Q\left( \sqrt{2},\sqrt{3} \right)\subseteq L$.
  \end{enumerate}
\end{solution}
\begin{problem}[Problem 4]
  Let $S = \set{n_1,\dots,n_k}$ be a finite set of positive integers greater than or equal to $2$, and let $K = \Q\left( \sqrt{n_1},\dots,\sqrt{n_k} \right)$.
  \begin{enumerate}[(a)]
    \item Prove that $\left[ K:\Q \right] = 2^{m}$ for some $m\leq k$, and the set 
      \begin{align*}
        P(S) &= \set{1}\cup\set{\sqrt{n} | n\text{ is a product of distinct elements of }S}
      \end{align*}
      spans $K$ over $\Q$.
    \item For each $0\leq j \leq k$, let $\Q_j = \Q\left( \sqrt{n_1},\dots,\sqrt{n_j} \right)$, setting $\Q_0 = \Q$. Prove that $\left[ K:\Q \right] < 2^{k}$ if and only if $n_1$ is a perfect square or there exists $2\leq i \leq k$ such that $\sqrt{n_i} = a + b\sqrt{n_{i-1}}$ for some $a,b\in \Q_{i-2}$.
    \item Assume that the elements of $S$ are pairwise relatively prime and none of them is a complete square. Prove that $\left[ K:\Q \right] = 2^{k}$.
  \end{enumerate}
\end{problem}
\begin{solution}\hfill
  \begin{enumerate}[(a)]
    \item Define $\Q_j = \Q\left( \sqrt{n_1},\dots,\sqrt{n_j} \right)$, where $\Q_0 = \Q$, as we do in part (b). Then, we have
      \begin{align*}
        \left[ K:\Q \right] &= \prod_{j=1}^{k} \left[ \Q_j:\Q_{j-1} \right].
      \end{align*}
      Each index value $\left[ \Q_j:\Q_{j-1} \right]$ is equal to either $1$ or $2$, meaning that $\left[ K:\Q \right] = 2^{m}$ for some $m\leq k$.

      To show that $P(S)$ spans $K$ over $\Q$, we start by observing that, by the definition of a field extension, we have
      \begin{align*}
        K &= \Q\left[ \sqrt{n_1},\dots,\sqrt{n_k} \right],
      \end{align*}
      meaning that $K$ is spanned by mixed monomials of the form $ax_1^{r_1}\cdots x_k^{r_k}$, where $a\in \Q$, evaluated at $\sqrt{n_1},\dots,\sqrt{n_k}$. Since $\left( \sqrt{n_k} \right)^2= n_k\in \Q$, we have that any mixed monomial $a\left( \sqrt{n_1} \right)^{r_1}\cdots \left( \sqrt{n_k} \right)^{r_k}$ may collapse any value of $r_i$ not equal to $0$ or $1$ into $a$, whence we have that $K$ is spanned by mixed monomials of the form
      \begin{align*}
        a\left( \sqrt{n_1} \right)^{\ve_1}\cdots \left( \sqrt{n_k} \right)^{\ve_k},
      \end{align*}
      where $\ve_i \in \set{0,1}$, which is exactly equal to $P(S)$.
    \item If $n_1$ is a perfect square, then $\Q\left( \sqrt{n_1} \right) = \Q$, whence $\left[ \Q_1:\Q \right] = 1$, meaning that $\left[ K:\Q \right] < 2^{k}$. Similarly, if there exists $2\leq i \leq k$ with $\sqrt{n_i} = a + b\sqrt{n_{i-1}}$ for some $a,b\in \Q_{i-2}$, then since $\Q_{i-2}\left( \sqrt{n_{i-1}} \right) = \Q_{i-1}$, there is a basis for $\Q_{i-1}$ over $\Q_{i-2}$ given by $\set{1,\sqrt{n_{i-1}}}$. If we have that $\sqrt{n_i} = a + b\sqrt{n_{i-1}}$ for some $a,b\in \Q_{i-2}$, then we have that 
      \begin{align*}
        \Q_{i} &= \Q_{i-1}\left( \sqrt{n_i} \right)\\
               &= \Q_{i-1},
      \end{align*}
      whence $\left[ \Q_{i}:\Q_{i-1} \right] = 1$, so we have yet again that $\left[ K:\Q \right] < 2^{k}$.

      Now, if $\left[ K:\Q \right] < 2^{k}$, then we write
      \begin{align*}
        \left[ K:\Q \right] &= \left[ \Q_k:\Q_{k-1} \right]\cdots \left[ \Q_1:\Q \right],
      \end{align*}
      and observe that each degree value is either $1$ or $2$. Since $\left[ K:\Q \right] < 2^{k}$, it follows that there is at least one factor that is equal to $1$. If $\left[ \Q_1:\Q \right] = 1$, then we have that $\sqrt{n_1}\in \Q$, meaning that $n_1$ is a perfect square. Otherwise, if there is some $i$ such that $\left[ \Q_i:\Q_{i-1} \right] = 1$, then since $\Q_{i-1} = \Q_{i-2}\left( \sqrt{n_{i-1}} \right)$, it follows that $\sqrt{n_i}\in \Q_{i-2}\left( \sqrt{n_{i-1}} \right)$, whence there are some $a,b\in \Q_{i-2}$ such that $\sqrt{n_i} = a + b\sqrt{n_{i-1}}$.
    \item Inductively, we start by assuming that $\left\vert S \right\vert = 0$, whence $K = \Q$ and so $\left[ K:\Q \right] = 1 = 2^{0}$.

      Consider now the tower of fields $\Q_0\subseteq \Q_1\subseteq\cdots\subseteq \Q_k$ as defined in part (b), and assume that each of $n_1,\dots,n_{k-1}$ are pairwise relatively prime integers that are not perfect squares. We will show that if $n_k$ is an integer that is not a perfect square and is coprime to all of $n_1,\dots,n_k$, then $\left[ \Q_k:\Q_{k-1} \right] = 2$. For if this weren't the case, then we would have $\left[ \Q_k:\Q \right] < 2^{k}$, meaning that from (b), we would have some $2\leq i \leq k$ such that $\sqrt{n_i} = a + b\sqrt{n_{i-1}}$ for some $a,b\in \Q_{i-2}$. This cannot be the case if $i = 2,\dots,k-1$ by the induction hypothesis, meaning that we would necessarily have $\sqrt{n_k} = a + b\sqrt{n_{k-1}}$ for some $a,b\in \Q_{k-2}$. Squaring gives
      \begin{align*}
        n_k &= a^2 + n_{k-1}b^2 + 2ab\sqrt{n_{k-1}}.
      \end{align*}
      If $a$ and $b$ are nonzero, then we have
      \begin{align*}
        \sqrt{n_{k-1}} &= \frac{n_k - a^2 - n_{k-1}b^2}{2ab}\\
                       &\in \Q_{k-2},
      \end{align*}
      which contradicts the induction hypothesis. If $a = 0$, then we have $n_k = n_{k-1}b^2$, which contradicts the assumption that $n_k$ is coprime to $n_1,\dots,n_{k-1}$, while if $b = 0$, then this contradicts the assumption that $n_k$ is not a perfect square.
  \end{enumerate}
\end{solution}
\begin{problem}[Problem 5]
  Let $F$ be a field, and let $\alpha$ be an algebraic element of odd degree over $F$. Prove that $F\left( \alpha^2 \right) = F\left( \alpha \right)$.
\end{problem}
\begin{solution}
  We have that $F\left( \alpha^2 \right)\subseteq F\left( \alpha \right)$. Let
  \begin{align*}
    \mu(x) &= x^{2k+1} + \chi(x)
  \end{align*}
  be the minimal polynomial for $\alpha$ over $F$, where $\chi(x)$ denotes lower order terms (which is not the zero polynomial since $\mu(x)$ is irreducible). Observe that we may express
  \begin{align*}
    \alpha^{2k+1} &= -\chi\left(\alpha\right).
  \end{align*}
  A basis for $F\left( \alpha \right)$ is given by
  \begin{align*}
    \set{1,\alpha,\dots,\alpha^{k},\alpha^{k+1},\dots,\alpha^{2k}},
  \end{align*}
  while the corresponding basis for $F\left( \alpha^2 \right)$ is given by
  \begin{align*}
    \set{1,\alpha^{2},\alpha^{4},\dots,\alpha^{2k},\alpha^{2k+2},\dots,\alpha^{4k}}.
  \end{align*}
  Yet, we observe that we may express $\set{\alpha^{2k+2},\dots,\alpha^{4k}}$ as $F$-linear combinations of lower order terms via the reduction process that $\alpha^{2k+1} = -\chi\left( \alpha \right)$ and divisions by $\alpha^{2\ell}$ with $0\leq \ell \leq 2k$ when necessary. Via this reduction process, we recover the $F$-basis for $F\left( \alpha \right)$, meaning that $F\left( \alpha^2 \right) = F\left( \alpha \right)$.
\end{solution}
\begin{problem}[Problem 8]
  Let $K/F$ be a field extension, and let $K_1,K_2$ be subfields of $K$ containing $F$ such that the extensions $K_1/F$ and $K_2/F$ are normal. Prove that the extensions $K_1K_2/F$ and $K_1\cap K_2/F$ are normal. Here, $K_1K_2$ is the composite field of $K_1$ and $K_2$.
\end{problem}
\begin{solution}
  Let $\alpha\in K_1\cap K_2$. Then, $\mu_{\alpha,F}(x)$ splits over $K_1$ by normality, and similarly it splits over $K_2$. Therefore, $\mu_{\alpha,F}(x)$ splits over $K_1\cap K_2$, so $K_1\cap K_2/F$ is normal.

  Since $K_1$ is normal, there is some subset $\Omega_1\subseteq F[x]$ of polynomials such that $K_1$ is the splitting field of this subset. Similarly, since $K_2$ is normal, there is $\Omega_2$ such that $K_2$ is the splitting field of this subset. Since $K_1K_2$ is the smallest subfield containing $K_1$ and $K_2$, it is thus the splitting field of $\Omega_1\cup \Omega_2$, whence $K_1K_2$ is normal.
\end{solution}
\end{document}
